\section{Introduction}
\label{sec:intro}

Fully test-time adaptation updates a pretrained model on unlabelled test batches as they arrive, for
example by minimizing prediction entropy over normalization parameters~\citep{wang2021tent}. Because every
update depends on the parameters left by the previous one, the final model is a function of the order of
the stream. Order dependence is not new: temporally correlated or non-i.i.d.\ streams degrade several TTA
methods~\citep{gong2022note,yuan2023rotta,niu2023sar,zhao2023pitfalls}, reordering alone changes results
in continual settings~\citep{wang2022cotta,dobler2023rmt,gao2023dda}, and resetting schemes are used to
contain long-horizon drift~\citep{press2023rdumb,niu2023sar,lim2026asr,hoang2024petta}. For plain SGD the
leading order-dependent term is also known: two consecutive steps on losses $L_a$ and $L_b$ differ from the
reversed order by $\eta^2(H_b g_a - H_a g_b)$ up to $O(\eta^3)$, a consequence of Taylor expansion that
appears in the analyses of Reptile, of implicit regularization in SGD and of data-ordering
attacks~\citep{nichol2018reptile,smith2021origin,shumailov2021ordering}, and more recently as a predictor
of curriculum order~\citep{sweeney2026lie}.

A natural idea is therefore to \emph{penalize} that term: restrict adaptation to directions along which
the local order difference, measured on model outputs, is small. Whether such a penalty is useful is an
empirical question, and it is easy to answer it wrongly. A restricted subspace moves less, uses extra
supervision if it is fitted on labelled data, and can change batch composition effects when orders are
formed by reshuffling samples. Any of these can reduce order sensitivity without the penalty playing a
role. This paper reports a controlled diagnostic of exactly this confound structure.

\paragraph{Question.} With the same meta supervision, the same adaptation budget and the same batch
membership, does an order penalty add utility over (i) a utility-only subspace that differs from the
candidate only by the penalty, (ii) simply using a smaller learning rate, and (iii) matching the size of
each update?

\paragraph{What we do.} We build a small synthetic TTA model in which every relevant quantity can be
controlled and audited (Section~\ref{sec:setting}): a linear softmax classifier with an adapted input affine
transform, a stream replayed in fixed orders whose batch membership is fixed so that only batch order
varies, a common terminal holdout for stationary streams, and current-regime holdouts plus
predict-then-update online error for genuine drift. The candidate and all controls are tuned on the same
development split with the same grid. The decision rule, seeds and minimum effects of interest were fixed
and committed before the run, and every reported number is exported programmatically from the raw run files.

\paragraph{What we find.} The pre-registered rule returned \emph{order penalty added utility not
supported} (Section~\ref{sec:experiments}). In two of three world seeds the penalized and the utility-only
fitters selected the same two pool directions and the same learning rate, which makes their trajectories
identical by construction; in the third seed the two differed and the candidate was worse. The
entropy-minimization baseline did not improve on no adaptation in this environment, which limits what the
environment can say about retaining adaptation gain.

\paragraph{Scope of the contribution.} We do not claim new theory: the order-difference expansion is
standard and we restate it only to define the candidate (Section~\ref{sec:analysis}). We do not claim
evidence about real-image TTA; no real data were used. The contribution is a controlled protocol, with
its controls and audits, and a negative diagnostic result reported together with the implementation-level
reason for it.
